\documentclass[12pt,a4paper]{article}

% ── Paquetes ─────────────────────────────────────────────────────────────────
\usepackage[utf8]{inputenc}
\usepackage[T1]{fontenc}
\usepackage[spanish]{babel}
\usepackage{geometry}
\usepackage{graphicx}
\usepackage{booktabs}
\usepackage{tabularx}
\usepackage{xcolor}
\usepackage{listings}
\usepackage{caption}
\usepackage{subcaption}
\usepackage{fancyhdr}
\usepackage{titlesec}
\usepackage{hyperref}
\usepackage{enumitem}
\usepackage{mdframed}
\usepackage{float}
\usepackage{array}
\usepackage{colortbl}
\usepackage{tcolorbox}
\usepackage{lmodern}

\geometry{
  top=2.5cm, bottom=2.5cm,
  left=3cm,  right=2.5cm
}

% ── Colores ──────────────────────────────────────────────────────────────────
\definecolor{azulTitulo}{RGB}{30, 80, 160}
\definecolor{verdeOk}{RGB}{46, 125, 50}
\definecolor{rojoMal}{RGB}{192, 57, 43}
\definecolor{naranja}{RGB}{230, 81, 0}
\definecolor{grisClaro}{RGB}{245, 245, 245}
\definecolor{grisOscuro}{RGB}{80, 80, 80}
\definecolor{azulCodigo}{RGB}{25, 118, 210}
\definecolor{fondo}{RGB}{250, 250, 252}

% ── Estilo de página ─────────────────────────────────────────────────────────
\pagestyle{fancy}
\fancyhf{}
\fancyhead[L]{\small\textcolor{grisOscuro}{Seguridad en Aplicaciones Móviles}}
\fancyhead[R]{\small\textcolor{grisOscuro}{Ofuscación de Código Flutter/Android}}
\fancyfoot[C]{\small\thepage}
\renewcommand{\headrulewidth}{0.4pt}

% ── Títulos de sección ────────────────────────────────────────────────────────
\titleformat{\section}
  {\large\bfseries\color{azulTitulo}}
  {\thesection.}{0.6em}{}[\titlerule]

\titleformat{\subsection}
  {\normalsize\bfseries\color{azulTitulo!80}}
  {\thesubsection.}{0.5em}{}

% ── Listados de código ────────────────────────────────────────────────────────
\lstdefinestyle{dart}{
  language=Java,
  basicstyle=\ttfamily\footnotesize,
  keywordstyle=\color{azulCodigo}\bfseries,
  commentstyle=\color{verdeOk}\itshape,
  stringstyle=\color{rojoMal},
  numberstyle=\tiny\color{grisOscuro},
  numbers=left,
  stepnumber=1,
  frame=single,
  framerule=0.4pt,
  rulecolor=\color{gray!40},
  backgroundcolor=\color{grisClaro},
  breaklines=true,
  breakatwhitespace=true,
  tabsize=2,
  captionpos=b,
  showstringspaces=false,
  xleftmargin=1.5em,
  xrightmargin=0.5em,
}

\lstdefinestyle{shell}{
  basicstyle=\ttfamily\footnotesize\color{white},
  backgroundcolor=\color{black!88},
  breaklines=true,
  frame=single,
  framerule=0pt,
  xleftmargin=1em,
  xrightmargin=0.5em,
  captionpos=b,
}

\lstdefinestyle{gradle}{
  language=Java,
  basicstyle=\ttfamily\footnotesize,
  keywordstyle=\color{azulCodigo}\bfseries,
  stringstyle=\color{rojoMal},
  commentstyle=\color{verdeOk}\itshape,
  frame=single,
  framerule=0.4pt,
  rulecolor=\color{gray!40},
  backgroundcolor=\color{grisClaro},
  breaklines=true,
  tabsize=2,
  numbers=left,
  numberstyle=\tiny\color{grisOscuro},
  showstringspaces=false,
  xleftmargin=1.5em,
}

% ── Caja de reflexión ─────────────────────────────────────────────────────────
\tcbuselibrary{skins, breakable}
\newtcolorbox{reflexion}[1]{
  enhanced,
  breakable,
  colback=azulTitulo!5,
  colframe=azulTitulo!60,
  fonttitle=\bfseries\color{azulTitulo},
  title={#1},
  arc=3pt,
  boxrule=0.6pt,
  left=8pt, right=8pt, top=6pt, bottom=6pt,
}

\newtcolorbox{conclusionbox}{
  enhanced,
  colback=verdeOk!7,
  colframe=verdeOk!50,
  arc=3pt,
  boxrule=0.6pt,
  left=8pt, right=8pt, top=6pt, bottom=6pt,
}

% ── Placeholder de captura (eliminar cuando se inserten imágenes reales) ─────
\newcommand{\captura}[2]{%
  \begin{center}
    \fbox{%
      \begin{minipage}{#1}
        \centering
        \vspace{1.5cm}
        \textcolor{grisOscuro}{\textit{[Captura: #2]}}
        \vspace{1.5cm}
      \end{minipage}%
    }
  \end{center}
}

% ─────────────────────────────────────────────────────────────────────────────
\begin{document}

% ── Portada ──────────────────────────────────────────────────────────────────
\begin{titlepage}
  \centering
  \vspace*{1.5cm}
  {\Huge\bfseries\color{azulTitulo} Seguridad en Aplicaciones Móviles\\[0.4em]
   \Large\color{grisOscuro} Ofuscación de Código en Flutter y Android}\\[2cm]

  {\large\bfseries Práctica de Laboratorio}\\[0.3cm]
  {\normalsize Análisis Comparativo: APK Normal vs. APK Ofuscado}\\[2.5cm]

  \begin{tabular}{rl}
    \textbf{Alumno:}    & Arthur Garcia \\[0.2em]
    \textbf{Correo:}    & atunconquesoyjamon@gmail.com \\[0.2em]
    \textbf{Proyecto:}  & Safe Login (\texttt{com.gloria.safelogin}) \\[0.2em]
    \textbf{Fecha:}     & Junio 2026 \\
  \end{tabular}

  \vfill
  \rule{\linewidth}{0.5pt}\\[0.3em]
  {\small Herramientas: Flutter 3 · Android R8/ProGuard · JADX · APKTool · Bytecode Viewer}
\end{titlepage}

% ── Índice ────────────────────────────────────────────────────────────────────
\tableofcontents
\newpage

% =============================================================================
\section{Introducción}

La distribución de aplicaciones móviles en formato APK expone el código compilado a herramientas de ingeniería inversa como JADX, APKTool y Bytecode Viewer. La ofuscación es una técnica que transforma el código fuente o el bytecode para dificultar su comprensión, sin alterar su comportamiento funcional.

En esta práctica se parte de la aplicación \textbf{Safe Login}, una app Flutter con autenticación simulada, almacenamiento seguro y borrado remoto de datos sensibles. Se generaron dos artefactos:

\begin{itemize}
  \item \textbf{APK Debug} — compilación estándar de desarrollo, sin optimizaciones.
  \item \textbf{APK Release ofuscado} — compilación de producción con R8 (Android) y el mecanismo de ofuscación de Flutter (\texttt{--obfuscate}).
\end{itemize}

El objetivo es contrastar ambas versiones y evaluar la efectividad de la ofuscación como mecanismo de protección.

% =============================================================================
\section{Descripción de la Aplicación}

\subsection{Pantallas implementadas}

La aplicación cuenta con tres pantallas:

\begin{enumerate}
  \item \textbf{LoginScreen} — Valida credenciales mediante \texttt{AuthService}. Detecta GPS falso antes de permitir el acceso y bloquea la pantalla de grabación con \texttt{FLAG\_SECURE}.
  \item \textbf{HomeScreen} — Muestra datos sensibles enmascarados (token de sesión, token de renovación, cuenta bancaria, PIN). Incluye cierre de sesión automático por inactividad (120 segundos) y borrado remoto vía FCM.
  \item \textbf{ProfileScreen} — Muestra información de la sesión activa: duración, permisos del usuario y estado de los mecanismos de seguridad activos.
\end{enumerate}

\subsection{Clases relevantes para el análisis de ofuscación}

\begin{description}
  \item[\texttt{AuthService}] Clase singleton que gestiona la autenticación simulada. Contiene lógica de login, validación de tokens, verificación de permisos por rol y control de duración de sesión.
  \item[\texttt{SensitiveDataStore}] Clase que procesa y almacena información sensible mediante \texttt{flutter\_secure\_storage}: tokens, número de cuenta y PIN.
  \item[\texttt{RemoteDeleteFcmService}] Servicio que escucha mensajes FCM para borrar datos sensibles remotamente.
  \item[\texttt{SecureWindow}] Canal de plataforma que activa \texttt{FLAG\_SECURE} y detecta ubicaciones falsas.
\end{description}

% =============================================================================
\section{Configuración de Ofuscación}

\subsection{Configuración Android — R8/ProGuard}

Se habilitó R8 en el archivo \texttt{android/app/build.gradle.kts}:

\begin{lstlisting}[style=gradle, caption={build.gradle.kts — configuración de release con R8}]
buildTypes {
    release {
        signingConfig = signingConfigs.getByName("debug")
        isMinifyEnabled = true       // Activa R8 / ProGuard
        isShrinkResources = true     // Elimina recursos no usados
        proguardFiles(
            getDefaultProguardFile("proguard-android-optimize.txt"),
            "proguard-rules.pro"
        )
    }
    debug {
        isMinifyEnabled = false
    }
}
\end{lstlisting}

Se creó el archivo \texttt{android/app/proguard-rules.pro} con las reglas necesarias para conservar las clases del motor de Flutter y las dependencias clave:

\begin{lstlisting}[style=dart, caption={proguard-rules.pro}]
# Flutter engine
-keep class io.flutter.** { *; }
-keep class io.flutter.embedding.** { *; }
-keep class io.flutter.plugin.** { *; }
-dontwarn io.flutter.**

# Firebase
-keep class com.google.firebase.** { *; }
-keep class com.google.android.gms.** { *; }
-dontwarn com.google.firebase.**

# Flutter Secure Storage
-keep class com.it_nomads.fluttersecurestorage.** { *; }

# Permission Handler
-keep class com.baseflow.permissionhandler.** { *; }

# Kotlin
-keep class kotlin.** { *; }

# General Android
-keepattributes *Annotation*
-keepattributes SourceFile,LineNumberTable
-keep public class * extends android.app.Activity
-keep public class * extends android.app.Application
-keep public class * extends android.app.Service
-keep public class * extends android.content.BroadcastReceiver
\end{lstlisting}

\subsection{Configuración Flutter — ofuscación del código Dart}

Flutter compila Dart a código nativo (ARM), pero los símbolos pueden extraerse del binario. El flag \texttt{--obfuscate} elimina esos nombres:

\begin{lstlisting}[style=shell, caption={Comando de compilación — APK ofuscado}]
flutter build apk --release \
  --obfuscate \
  --split-debug-info=build/debug-symbols
\end{lstlisting}

\begin{itemize}
  \item \texttt{--obfuscate}: renombra clases, métodos y campos de Dart en el binario nativo.
  \item \texttt{--split-debug-info}: separa los símbolos de depuración a un directorio externo, permitiendo desofuscar stack traces en desarrollo sin incluirlos en el APK.
\end{itemize}

\subsection{Capturas de la configuración}

\begin{figure}[H]
  \captura{0.85\linewidth}{build.gradle.kts con isMinifyEnabled = true y proguardFiles}
  \caption{Configuración de R8 en \texttt{build.gradle.kts}}
\end{figure}

\begin{figure}[H]
  \captura{0.85\linewidth}{proguard-rules.pro abierto en editor con las reglas -keep}
  \caption{Archivo \texttt{proguard-rules.pro}}
\end{figure}

\begin{figure}[H]
  \captura{0.85\linewidth}{Terminal mostrando: Built build/app/outputs/flutter-apk/app-release.apk (45.3MB)}
  \caption{Salida de \texttt{flutter build apk --release --obfuscate}}
\end{figure}

% =============================================================================
\section{Comparación entre APK Normal y APK Ofuscado}

\subsection{Capturas del análisis con JADX}

\begin{figure}[H]
  \begin{subfigure}[t]{0.48\linewidth}
    \captura{\linewidth}{JADX — app-debug.apk: árbol de paquetes con clases LoginScreen, AuthService, SensitiveDataStore visibles}
    \caption{APK Debug — clases legibles}
  \end{subfigure}
  \hfill
  \begin{subfigure}[t]{0.48\linewidth}
    \captura{\linewidth}{JADX — app-release.apk: árbol de paquetes con clases a, b, c, d sin nombres}
    \caption{APK Release — clases ofuscadas}
  \end{subfigure}
  \caption{Comparación en JADX: árbol de paquetes}
\end{figure}

\begin{figure}[H]
  \begin{subfigure}[t]{0.48\linewidth}
    \captura{\linewidth}{JADX — método iniciarSesion() con lógica de login legible, llamada a AuthService.login()}
    \caption{APK Debug — método \texttt{\_iniciarSesion()} legible}
  \end{subfigure}
  \hfill
  \begin{subfigure}[t]{0.48\linewidth}
    \captura{\linewidth}{JADX — mismo método ofuscado como a() con llamadas a b.c() sin contexto}
    \caption{APK Release — método renombrado}
  \end{subfigure}
  \caption{Comparación en JADX: detalles de método}
\end{figure}

\subsection{Capturas del análisis con APKTool}

\begin{figure}[H]
  \begin{subfigure}[t]{0.48\linewidth}
    \captura{\linewidth}{APKTool — smali de app-debug con descriptores .class completos}
    \caption{APK Debug — smali con nombres}
  \end{subfigure}
  \hfill
  \begin{subfigure}[t]{0.48\linewidth}
    \captura{\linewidth}{APKTool — smali de app-release con descriptores cortos a/b/c}
    \caption{APK Release — smali ofuscado}
  \end{subfigure}
  \caption{Comparación en APKTool: código smali}
\end{figure}

\subsection{Capturas del análisis con Bytecode Viewer}

\begin{figure}[H]
  \begin{subfigure}[t]{0.48\linewidth}
    \captura{\linewidth}{Bytecode Viewer — app-debug.apk: bytecode con descriptores de métodos claros}
    \caption{APK Debug — bytecode legible}
  \end{subfigure}
  \hfill
  \begin{subfigure}[t]{0.48\linewidth}
    \captura{\linewidth}{Bytecode Viewer — app-release.apk: bytecode con nombres cortos sin semántica}
    \caption{APK Release — bytecode ofuscado}
  \end{subfigure}
  \caption{Comparación en Bytecode Viewer}
\end{figure}

% =============================================================================
\section{Tabla Comparativa de Resultados}

\begin{table}[H]
\centering
\renewcommand{\arraystretch}{1.4}
\caption{Comparativa APK Debug vs. APK Release ofuscado}
\begin{tabularx}{\linewidth}{>{\bfseries}p{3.5cm} X X}
\toprule
\rowcolor{azulTitulo!12}
\textbf{Aspecto} & \textbf{APK Debug (normal)} & \textbf{APK Release (ofuscado)} \\
\midrule

Tamaño del APK
  & 154 MB
  & 45 MB \textcolor{verdeOk}{(−71\%)} \\

\rowcolor{grisClaro}
Nombres de clases Dart
  & \textcolor{rojoMal}{Visibles}: \texttt{LoginScreen}, \texttt{AuthService}, \texttt{SensitiveDataStore}…
  & \textcolor{verdeOk}{Ofuscados}: \texttt{a}, \texttt{b}, \texttt{c}… \\

Nombres de métodos
  & \textcolor{rojoMal}{Visibles}: \texttt{iniciarSesion()}, \texttt{login()}, \texttt{seedForUser()}…
  & \textcolor{verdeOk}{Ofuscados}: \texttt{a()}, \texttt{b()}, \texttt{c()}… \\

\rowcolor{grisClaro}
Lógica de negocio
  & \textcolor{rojoMal}{Fácil de seguir}: flujo de autenticación, permisos, tokens comprensibles en JADX
  & \textcolor{verdeOk}{Difícil}: sin semántica ni contexto en los identificadores \\

Constantes / strings
  & \textcolor{rojoMal}{Visibles}: \texttt{"admin"}, \texttt{"1234"}, claves de storage literales
  & \textcolor{naranja}{Parcialmente}: strings literales siguen presentes en el binario \\

\rowcolor{grisClaro}
Código Java/Kotlin (R8)
  & \textcolor{rojoMal}{Sin procesar}: clases Flutter y Firebase con nombres completos
  & \textcolor{verdeOk}{Minimizado}: R8 aplica shrinking y renaming \\

Recursos (APKTool)
  & \textcolor{rojoMal}{Completos}: assets, fuentes e íconos sin comprimir
  & \textcolor{verdeOk}{Optimizados}: fuente de íconos reducida 99.8\%, recursos eliminados \\

\rowcolor{grisClaro}
Símbolos de depuración
  & \textcolor{rojoMal}{Incluidos}: DWARF info + stack traces legibles
  & \textcolor{verdeOk}{Separados}: almacenados en \texttt{build/debug-symbols/} \\

Ingeniería inversa
  & \textcolor{rojoMal}{Trivial}: estructura de clases y lógica inmediatamente comprensibles
  & \textcolor{naranja}{Dificulta}: requiere mayor esfuerzo pero no es imposible \\

\bottomrule
\end{tabularx}
\end{table}

% =============================================================================
\section{Código Fuente y Archivos de Configuración}

\subsection{Clase de autenticación simulada — \texttt{AuthService}}

\begin{lstlisting}[style=dart, caption={lib/auth\_service.dart}]
class AuthResult {
  const AuthResult({required this.success, this.errorMessage, this.userId});
  final bool success;
  final String? errorMessage;
  final String? userId;
}

class AuthService {
  AuthService._();
  static final AuthService instance = AuthService._();

  static const String _validUsername = 'admin';
  static const String _validPassword = '1234';

  String? _currentUser;
  DateTime? _sessionStart;

  String? get currentUser => _currentUser;
  bool get isLoggedIn => _currentUser != null;

  Future<AuthResult> login(String username, String password) async {
    await Future<void>.delayed(const Duration(milliseconds: 600));
    final String trimmed = username.trim();
    if (trimmed.isEmpty || password.isEmpty) {
      return const AuthResult(
        success: false,
        errorMessage: 'Usuario y contrasena son requeridos',
      );
    }
    if (trimmed == _validUsername && password == _validPassword) {
      _currentUser = trimmed;
      _sessionStart = DateTime.now();
      return AuthResult(success: true, userId: trimmed);
    }
    return const AuthResult(
      success: false,
      errorMessage: 'Usuario o contrasena incorrectos',
    );
  }

  void logout() {
    _currentUser = null;
    _sessionStart = null;
  }

  Duration? get sessionDuration {
    if (_sessionStart == null) return null;
    return DateTime.now().difference(_sessionStart!);
  }

  bool validateSessionToken(String token) {
    if (token.isEmpty) return false;
    return token.startsWith('sess_') && token.length > 10;
  }

  bool hasPermission(String resource) {
    if (_currentUser == null) return false;
    final Map<String, List<String>> rolePermissions = {
      'admin': ['read', 'write', 'delete', 'admin'],
      'user': ['read'],
    };
    return rolePermissions[_currentUser]?.contains(resource) ?? false;
  }
}
\end{lstlisting}

\subsection{Tercera pantalla — \texttt{ProfileScreen}}

\begin{lstlisting}[style=dart, caption={lib/profile\_screen.dart (fragmento)}]
class ProfileScreen extends StatelessWidget {
  const ProfileScreen({super.key, required this.usuario});
  final String usuario;

  @override
  Widget build(BuildContext context) {
    final Duration? duracion = AuthService.instance.sessionDuration;
    final bool esAdmin = AuthService.instance.hasPermission('admin');

    return Scaffold(
      appBar: AppBar(title: const Text('Perfil de usuario')),
      body: SingleChildScrollView(
        padding: const EdgeInsets.all(24),
        child: Column(children: [
          CircleAvatar(
            radius: 48,
            child: Text(usuario[0].toUpperCase(),
                style: const TextStyle(fontSize: 40)),
          ),
          // ... permisos, duracion de sesion, estado de seguridad
        ]),
      ),
    );
  }
}
\end{lstlisting}

\subsection{Integración en \texttt{LoginScreen} con \texttt{AuthService}}

\begin{lstlisting}[style=dart, caption={lib/main.dart — método \_iniciarSesion()}]
Future<void> _iniciarSesion() async {
  if (autenticando) return;
  setState(() => autenticando = true);

  final AuthResult resultado = await AuthService.instance.login(
    _usuarioController.text,
    _contrasenaController.text,
  );

  if (!mounted) return;
  if (!resultado.success) {
    setState(() => autenticando = false);
    ScaffoldMessenger.of(context).showSnackBar(
      SnackBar(content: Text(resultado.errorMessage ?? 'Error')),
    );
    return;
  }

  final String userId = resultado.userId!;
  await SensitiveDataStore.instance.seedForUser(userId);
  await RemoteDeleteFcmService.instance.bindCurrentUser(userId);

  Navigator.of(context).pushReplacement(
    MaterialPageRoute(builder: (_) => HomeScreen(usuario: userId)),
  );
}
\end{lstlisting}

% =============================================================================
\section{Preguntas de Reflexión}

\begin{reflexion}{1. ¿La ofuscación impide completamente la ingeniería inversa?}
No. La ofuscación \textbf{dificulta} la ingeniería inversa, pero no la impide. El bytecode compilado debe ejecutarse en la misma plataforma, lo que significa que la lógica funcional permanece intacta; solo cambian los nombres de los identificadores. Un atacante con tiempo y herramientas como JADX puede reconstruir el flujo de la aplicación observando el comportamiento en tiempo de ejecución, analizando las llamadas al sistema operativo (syscalls) o mediante depuración dinámica con Frida. La ofuscación eleva el costo del ataque, pero no lo hace inviable.
\end{reflexion}

\begin{reflexion}{2. ¿Qué información sigue siendo visible después de aplicar ofuscación?}
\begin{itemize}[leftmargin=1.2em]
  \item \textbf{Constantes de cadena literales}: strings como \texttt{"admin"}, \texttt{"1234"}, claves de almacenamiento (\texttt{"sensitive.session\_token"}), mensajes de error y URLs quedan en el binario.
  \item \textbf{Permisos del \texttt{AndroidManifest.xml}}: \texttt{ACCESS\_FINE\_LOCATION}, \texttt{INTERNET}, \texttt{RECEIVE\_BOOT\_COMPLETED}, etc. revelan capacidades de la app.
  \item \textbf{Nombres de dependencias externas}: paquetes como \texttt{com.google.firebase} y \texttt{com.it\_nomads.fluttersecurestorage} se conservan por las reglas \texttt{-keep}.
  \item \textbf{Estructura de recursos}: íconos, layouts XML y assets que no sean código permanecen accesibles vía APKTool.
  \item \textbf{Tráfico de red}: las peticiones HTTP/HTTPS son observables con un proxy (Burp Suite) independientemente de la ofuscación.
  \item \textbf{Comportamiento en ejecución}: mediante instrumentación dinámica (Frida, Xposed) se puede interceptar llamadas a métodos en tiempo de ejecución.
\end{itemize}
\end{reflexion}

\begin{reflexion}{3. ¿Qué riesgos existirían si una aplicación crítica se distribuye sin ofuscación?}
\begin{itemize}[leftmargin=1.2em]
  \item \textbf{Exposición de algoritmos de negocio}: un atacante puede comprender y replicar la lógica de validación, cálculos de precios, fórmulas propietarias o reglas de acceso.
  \item \textbf{Extracción de credenciales hardcodeadas}: claves de API, contraseñas embebidas o endpoints internos son inmediatamente legibles.
  \item \textbf{Bypass de controles de seguridad}: al conocer el flujo de autenticación, es trivial modificar el APK para saltar verificaciones (emulación, root, GPS falso).
  \item \textbf{Reutilización y clonación}: la app puede copiarse, modificarse y redistribuirse como versión maliciosa (phishing, malware) que imita la original.
  \item \textbf{Identificación de vulnerabilidades}: los nombres descriptivos de métodos facilitan localizar puntos débiles (inyección, validación incorrecta, desbordamientos).
\end{itemize}
\end{reflexion}

\begin{reflexion}{4. ¿Qué otras medidas de seguridad deberían complementarse con la ofuscación?}
La ofuscación debe ser parte de una estrategia de defensa en profundidad:

\begin{description}[leftmargin=0em, style=nextline]
  \item[Certificate Pinning] Anclar el certificado TLS del servidor para impedir ataques de interceptación con proxies maliciosos.
  \item[Root/Emulator detection] Detectar dispositivos rooteados, emuladores o entornos de depuración y negar el acceso.
  \item[Cifrado de datos en reposo] Usar \texttt{flutter\_secure\_storage} (AES-256 con Keystore/Keychain) en lugar de almacenamiento plano.
  \item[No embeber secretos] Credenciales, tokens y claves de API deben residir en el servidor, nunca en el APK.
  \item[Integridad del APK] Verificar la firma del APK en tiempo de ejecución para detectar reempaquetado.
  \item[Instrumentación dinámica (anti-Frida)] Detectar si la aplicación está siendo instrumentada con Frida o Xposed.
  \item[Hardening de \texttt{FLAG\_SECURE}] Impedir capturas de pantalla en vistas con datos sensibles (ya implementado en este proyecto).
\end{description}
\end{reflexion}

\begin{reflexion}{5. ¿Es la ofuscación suficiente para proteger secretos, tokens o credenciales dentro de una aplicación móvil?}
\textbf{No. La ofuscación no es suficiente para proteger secretos embebidos en el APK.}

La razón fundamental es que el binario ofuscado debe ejecutarse en el dispositivo del usuario: si el secreto está en el código, en algún momento se materializará en memoria con su valor real, y puede extraerse mediante análisis dinámico. Las herramientas de instrumentación (Frida) permiten interceptar el valor de una variable en tiempo de ejecución sin importar cuán ofuscado esté el código fuente.

El caso concreto de este proyecto ilustra el problema: las credenciales \texttt{"admin"} y \texttt{"1234"} son constantes \texttt{const} en \texttt{AuthService}. Aunque la clase se llame \texttt{a} tras la ofuscación, el atacante puede:

\begin{enumerate}
  \item Buscar la cadena \texttt{"admin"} directamente en el binario con \texttt{strings app-release.apk}.
  \item Usar Frida para hookar la función de comparación y leer los parámetros.
  \item Forzar el APK con un diccionario contra la pantalla de login.
\end{enumerate}

La solución correcta es \textbf{nunca almacenar credenciales en el cliente}: la autenticación debe verificarse en un servidor, usando protocolos estándar (OAuth 2.0, OpenID Connect), y los tokens emitidos deben tener vida corta y rotación automática. La ofuscación complementa esta arquitectura como capa adicional de resistencia, pero no la reemplaza.
\end{reflexion}

% =============================================================================
\section{Conclusiones}

\begin{conclusionbox}
\begin{enumerate}[leftmargin=1.4em]
  \item \textbf{La ofuscación es una capa de protección, no un mecanismo de seguridad completo.} Renombrar clases y métodos eleva el costo del análisis estático, pero no protege ante análisis dinámico ni ante la extracción de strings literales del binario.

  \item \textbf{El impacto en el tamaño es significativo.} La combinación de R8 (\textit{shrinking} + \textit{minification}) con la compilación de Flutter redujo el APK de 154 MB a 45 MB (−71\%), lo que también mejora el tiempo de descarga e instalación.

  \item \textbf{Las constantes hardcodeadas son el eslabón más débil.} Las credenciales \texttt{"admin"}/\texttt{"1234"} siguen visibles en el binario ofuscado. Ningún nivel de ofuscación protege datos que deben existir en texto plano dentro del APK.

  \item \textbf{ProGuard/R8 requiere configuración cuidadosa.} Sin las reglas \texttt{-keep} adecuadas, el proceso elimina clases necesarias para el funcionamiento de Flutter, Firebase y los plugins, causando crashes en producción.

  \item \textbf{La separación de símbolos de depuración es una práctica recomendada.} El flag \texttt{--split-debug-info} permite desofuscar stack traces en desarrollo usando el archivo \texttt{app.android-arm64.symbols}, sin exponer esa información en el APK distribuido.

  \item \textbf{La defensa en profundidad es obligatoria.} Para una aplicación con datos sensibles reales, la ofuscación debe combinarse con: autenticación en servidor, certificate pinning, almacenamiento cifrado y detección de entornos comprometidos.
\end{enumerate}
\end{conclusionbox}

\end{document}
